\documentclass[conference]{IEEEtran}
\IEEEoverridecommandlockouts
\usepackage{cite}
\usepackage{amsmath,amssymb,amsfonts}
\usepackage{graphicx}
\usepackage{textcomp}
\usepackage{xcolor}
\usepackage[T1]{fontenc}
\usepackage[utf8]{inputenc}
\usepackage[brazil]{babel}
\usepackage{hyperref}
\usepackage{booktabs}

\def\BibTeX{{\rm B\kern-.05em{\sc i\kern-.025em b}\kern-.08em
    T\kern-.1667em\lower.7ex\hbox{E}\kern-.125emX}}

\begin{document}

% ─────────────────────────────────────────────────────────────
\title{GraphRAG para Consultas sobre Eventos Climáticos Extremos
no Brasil: uma Avaliação Comparativa}

\author{
\IEEEauthorblockN{Julio Rafael}
\IEEEauthorblockA{
  \textit{Ciência da Computação} \\
  \textit{Universidade Tiradentes (UNIT)} \\
  Aracaju, Brasil \\
  juliorafaelnft@hotmail.com
}
}

\maketitle

% ─────────────────────────────────────────────────────────────
\begin{abstract}
O acesso eficiente a registros históricos de eventos climáticos
extremos é essencial para gestão de riscos e políticas públicas.
Modelos de Linguagem de Grande Escala (LLMs) possibilitam consultas
em linguagem natural sobre esses dados, mas carecem de mecanismos
para recuperar fatos específicos de municípios ou períodos.
Este trabalho propõe um pipeline GraphRAG que estrutura dados do
INMET e do CEMADEN em um grafo de conhecimento Neo4j e os expõe
via interface conversacional.
Comparamos três estratégias: LLM sem recuperação (No~RAG),
RAG vetorial (Vector~RAG) e o pipeline proposto (Graph~RAG),
avaliados com o framework RAGAS em cinco perguntas sobre
eventos climáticos no Brasil.
Os resultados indicam que o GraphRAG supera os demais em
fidelidade e abrangência de contexto, com ganho médio de
\textbf{XX~p.p.} em \textit{faithfulness} sobre o Vector~RAG
e de \textbf{YY~p.p.} sobre o baseline sem recuperação.
\end{abstract}

\begin{IEEEkeywords}
GraphRAG, eventos climáticos extremos, grafo de conhecimento,
RAGAS, LLM, Brasil, Neo4j
\end{IEEEkeywords}

% ─────────────────────────────────────────────────────────────
\section{Introdução}

O Brasil está entre os países mais vulneráveis a eventos climáticos
extremos: enchentes, deslizamentos, secas e ondas de calor causam
perdas humanas e econômicas expressivas a cada ano~\cite{atkins2024generative}.
Agências nacionais como o Instituto Nacional de Meteorologia
(INMET) e o Centro Nacional de Monitoramento e Alertas de Desastres
Naturais (CEMADEN) produzem volumes crescentes de dados
estruturados e relatórios textuais sobre esses eventos.
Ainda assim, o acesso a essas informações exige conhecimento
técnico de SQL/Cypher ou leitura manual de boletins em PDF,
criando uma barreira para gestores, jornalistas e pesquisadores.

Modelos de Linguagem de Grande Escala (LLMs), como GPT-4o,
permitem consultas em linguagem natural, mas apresentam dois
problemas críticos quando aplicados isoladamente a dados
operacionais: (i)~\textbf{alucinação}, geração de fatos
plausíveis mas incorretos; e (ii)~\textbf{desatualização},
pois o conhecimento do modelo encerra-se em sua data de corte
de treinamento e não inclui dados regionais específicos.

Geração Aumentada por Recuperação (RAG) mitiga esses problemas
ao recuperar trechos relevantes de uma base externa antes de
gerar a resposta.
O RAG vetorial é a abordagem mais comum: textos são convertidos em
embeddings e armazenados em um índice, e os mais similares à consulta
são recuperados.
Contudo, perguntas relacionais --- \emph{``Quais municípios do RS
tiveram mais enchentes entre 2020 e 2024?''} --- exigem cruzamentos
entre entidades (municípios, eventos, fenômenos, datas) que índices
vetoriais planos não representam adequadamente.

Edge et al.~\cite{edge2024graphrag} propuseram o GraphRAG, que
organiza o conhecimento em um grafo de entidades e relações.
Durante a recuperação, o sistema percorre o grafo para reunir
subgrafos contextualmente ricos, permitindo respostas mais precisas
a consultas relacionais.
Até o momento, nenhum trabalho aplicou GraphRAG a dados climáticos
operacionais brasileiros.

\textbf{Contribuições deste trabalho:}
\begin{enumerate}
  \item Um grafo de conhecimento climático construído a partir de
        dados reais do INMET e boletins do CEMADEN, armazenado no
        Neo4j com esquema orientado a eventos extremos brasileiros.
  \item Um pipeline GraphRAG com extração de entidades via LLM,
        consultas Cypher parametrizadas e geração de resposta em
        streaming.
  \item Uma avaliação comparativa quantitativa com o framework
        RAGAS, cobrindo três estratégias (No~RAG, Vector~RAG,
        Graph~RAG) em cinco perguntas de domínio.
\end{enumerate}

O restante do artigo está organizado como segue:
a Seção~\ref{sec:relacionados} revisa trabalhos relacionados;
a Seção~\ref{sec:metodologia} descreve o pipeline proposto;
a Seção~\ref{sec:resultados} apresenta e discute os resultados;
e a Seção~\ref{sec:conclusao} conclui com direções futuras.

% ─────────────────────────────────────────────────────────────
\section{Trabalhos Relacionados}
\label{sec:relacionados}

A interseção entre Modelos de Linguagem de Grande Escala (LLMs)
e mudanças climáticas tem atraído crescente atenção da comunidade
científica, com esforços concentrados na melhoria da precisão,
redução de alucinações e ampliação do acesso a informações
climáticas.

Vaghefi et al.~\cite{vaghefi2023chatclimate} introduziram o
ChatClimate, uma IA conversacional fundamentada no Sexto Relatório
de Avaliação do IPCC (AR6) por meio de Geração Aumentada por
Recuperação (RAG). Comparando uma versão standalone, uma versão
híbrida e o GPT-4 puro em 13 perguntas elaboradas por especialistas,
os autores demonstraram que o enriquecimento de LLMs com memória
externa específica do domínio reduz significativamente alucinações
e respostas desatualizadas, com o modelo híbrido apresentando as
respostas mais precisas.

Atkins et al.~\cite{atkins2024generative} investigaram o potencial
de ferramentas de IA generativa de propósito geral, especificamente
GPT-3.5 e GPT-4, para o letramento climático.
Utilizando um conjunto de dados global com 191 países e validando
as respostas contra o índice de risco INFORM baseado no IPCC AR6,
o estudo encontrou precisão razoável para inundações e ciclones,
mas menor concordância para secas, evidenciando vieses persistentes
em favor de contextos anglófonos e ocidentais.

Juhasz et al.~\cite{juhasz2024responsible} propuseram um framework
RAG responsável para documentos de legislação e política climática,
introduzindo um conjunto de avaliação com dimensões de fidelidade,
alinhamento de política, formatação e resposta do sistema.
Avaliando múltiplos modelos de código aberto sobre 550 documentos
da base Climate Policy Radar, os autores defendem uma abordagem de
implantação em profundidade e uma interface que expõe pontuações
de avaliação e citações de fontes em tempo real.

Hsu et al.~\cite{hsu2024chatnetzero} apresentaram o ChatNetZero,
um chatbot baseado em RAG especializado em compromissos climáticos
de net-zero, incorporando módulos dedicados de anti-alucinação e
referência.
Ao comparar o ChatNetZero com GPT-4, Gemini, Coral e ChatClimate,
os autores constataram que, embora o ChatNetZero tenha alcançado
maior precisão factual, especialistas preferiram consistentemente
respostas mais longas e contextuais, revelando um viés de
verbosidade na avaliação humana de saídas de LLMs.

Enquanto os trabalhos acima se baseiam em RAG vetorial sobre
documentos climáticos, Edge et al.~\cite{edge2024graphrag}
introduziram o GraphRAG, um framework que estrutura o conhecimento
como um grafo de entidades em vez de um índice vetorial plano.
Ao percorrer relações do grafo durante a recuperação, o GraphRAG
permite respostas mais precisas a consultas relacionais --- como
identificar quais municípios foram afetados por eventos extremos
compostos --- tarefa em que o RAG convencional frequentemente falha.

A avaliação de sistemas RAG foi sistematizada por
Es et al.~\cite{es2024ragas}, que propuseram o RAGAS, um framework
automatizado que mede fidelidade, relevância da resposta e
abrangência do contexto sem exigir gabaritos anotados manualmente.
O RAGAS foi adotado em trabalhos subsequentes de IA climática
\cite{juhasz2024responsible} como benchmark padrão para comparação
de estratégias de recuperação.

Comparações diretas entre Vector~RAG e GraphRAG têm emergido
recentemente.
Ahmad et al.~\cite{ahmad2025benchmarking} avaliaram pipelines
Vector, Graph e Hybrid RAG no domínio de redes de acesso de rádio
(ORAN) usando as métricas RAGAS, reportando que o GraphRAG supera
o Vector~RAG em Faithfulness ($+$4 p.p.) e Factual Correctness
($+$2 p.p.), enquanto scores de Context Relevance permaneceram
baixos (0,04--0,11) em ambos os sistemas.
Han et al.~\cite{han2025ragvsgraphrag} realizaram uma avaliação
sistemática em benchmarks de QA (Natural Questions, HotpotQA),
encontrando diferenças modestas entre as abordagens ($<$2 p.p.
em F1), sugerindo que os ganhos do GraphRAG dependem fortemente
da natureza das consultas.
Zeng et al.~\cite{zeng2025unbiased} alertam que ganhos relatados
pelo GraphRAG sobre métodos de RAG convencional podem ser
significativamente menores após correção de vieses de avaliação.

Apesar desses avanços, nenhum trabalho anterior aplicou GraphRAG
a dados climáticos operacionais brasileiros.
Os sistemas existentes concentram-se em relatórios globais do IPCC
ou documentos de política de net-zero, deixando uma lacuna para
acesso estruturado e relacional aos registros produzidos por agências
nacionais como o INMET e o CEMADEN.
O presente trabalho endereça essa lacuna construindo um grafo de
conhecimento a partir de dados de estações meteorológicas e boletins
de eventos extremos, e avaliando um pipeline GraphRAG frente a uma
baseline de RAG vetorial e a um LLM isolado.

% ─────────────────────────────────────────────────────────────
\section{Metodologia}
\label{sec:metodologia}

\subsection{Fontes de Dados}

O sistema utiliza duas fontes de dados públicas brasileiras:

\textbf{INMET} --- O Instituto Nacional de Meteorologia disponibiliza
séries históricas horárias e diárias de estações meteorológicas
automáticas distribuídas por todos os estados.
Os dados incluem temperatura máxima/mínima/média, precipitação,
umidade relativa, velocidade do vento e pressão atmosférica.
Para este trabalho, foram coletados registros de estações das regiões
Sul, Sudeste e Centro-Oeste, cobrindo o período de 2019 a 2024.

\textbf{CEMADEN} --- O Centro Nacional de Monitoramento e Alertas
de Desastres Naturais publica boletins mensais em PDF descrevendo
eventos extremos (enchentes, deslizamentos, secas, granizo, vendavais
e outros) com informações de município, estado, data, número de
afetados e mortos.
Os boletins foram processados com \texttt{pdfplumber} para extração
de texto, e as informações estruturadas foram extraídas via
Claude claude-sonnet-4-6 (API Anthropic).

\subsection{Grafo de Conhecimento}

O conhecimento climático é representado em um grafo Neo4j com o
seguinte esquema de nós e relações (Fig.~\ref{fig:esquema_grafo}):

\begin{itemize}
  \item \textbf{Nós}: \texttt{Municipio}, \texttt{Estado},
        \texttt{Evento}, \texttt{Fenomeno}, \texttt{Estacao}
  \item \textbf{Relações}: \texttt{PERTENCE\_A}, \texttt{SOFREU},
        \texttt{CAUSADO\_POR}, \texttt{REGISTRADO\_EM}
\end{itemize}

Cada nó \texttt{Evento} armazena tipo, data, valor medido,
descrição e fonte; cada nó \texttt{Municipio} carrega nome,
código IBGE e sigla de estado.
Esse esquema permite consultas relacionais como:
\emph{``MATCH (m:Municipio)-[:SOFREU]->(e:Evento)-[:CAUSADO\_POR]
->(f:Fenomeno) WHERE f.chave = `enchente' RETURN m.nome, count(e)
ORDER BY count(e) DESC''}.

\begin{figure}[htbp]
\centerline{\includegraphics[width=0.45\textwidth]{fig_esquema_grafo.png}}
\caption{Esquema do grafo de conhecimento climático no Neo4j.}
\label{fig:esquema_grafo}
\end{figure}

\subsection{Pipeline GraphRAG}

O pipeline é implementado em TypeScript com Bun e exposto como
API HTTP via Hono. Dado uma pergunta em linguagem natural, o
pipeline executa quatro etapas (Fig.~\ref{fig:pipeline}):

\begin{enumerate}
  \item \textbf{Extração de entidades} --- O modelo GPT-4o-mini
        extrai entidades geográficas (municípios, estados),
        tipos de evento e períodos temporais usando
        \textit{structured output} com esquema Zod.
  \item \textbf{Recuperação do subgrafo} --- As entidades extraídas
        parametrizam consultas Cypher pré-definidas no Neo4j,
        retornando registros relevantes do grafo.
  \item \textbf{Construção do contexto} --- Os registros são
        formatados em texto estruturado que lista município, tipo
        de evento, data, valor medido, fenômeno e fonte.
  \item \textbf{Geração da resposta} --- GPT-4o recebe o contexto
        e a pergunta e gera a resposta em streaming via
        Server-Sent Events (SSE), com instrução explícita de não
        inventar dados ausentes no contexto.
\end{enumerate}

\begin{figure}[htbp]
\centerline{\includegraphics[width=0.45\textwidth]{fig_pipeline.png}}
\caption{Pipeline GraphRAG: extração de entidades, recuperação
  Cypher, construção de contexto e geração de resposta.}
\label{fig:pipeline}
\end{figure}

\subsection{Baseline Vector RAG}

Para o Vector~RAG, os eventos são indexados no Neo4j Vector Index
usando embeddings \texttt{text-embedding-3-small} (OpenAI) sobre
o campo \texttt{descricao} de cada nó \texttt{Evento}.
A recuperação usa o retriever LangChain com $k=10$ documentos
e uma \textit{RetrievalChain} com GPT-4o como gerador.
Diferentemente do GraphRAG, esta abordagem não percorre relações
do grafo; ela retorna os eventos textuais mais similares à consulta.

\subsection{Baseline No RAG}

O baseline sem recuperação consiste em enviar a pergunta diretamente
ao GPT-4o sem qualquer contexto externo, dependendo exclusivamente
do conhecimento paramétrico do modelo.
Esse cenário representa o comportamento de um LLM de propósito geral
não especializado em dados climáticos brasileiros.

\subsection{Avaliação com RAGAS}

A avaliação utiliza o framework RAGAS~\cite{es2024ragas} com
quatro métricas:

\begin{itemize}
  \item \textbf{Faithfulness} (fidelidade): mede se as afirmações
        da resposta estão sustentadas pelo contexto recuperado,
        detectando alucinações. Valor entre 0 e 1.
  \item \textbf{Answer Relevancy} (relevância da resposta): mede
        se a resposta aborda de fato a pergunta formulada, penalizando
        respostas incompletas ou tangenciais.
  \item \textbf{Context Recall} (abrangência do contexto): mede a
        fração do gabarito coberta pelo contexto recuperado,
        indicando a completude da recuperação.
  \item \textbf{Context Precision} (precisão do contexto): mede a
        proporção de chunks recuperados que são relevantes,
        penalizando ruído na recuperação.
\end{itemize}

Para o No~RAG, faithfulness e answer\_relevancy são calculados sem
contexto; context\_recall e context\_precision recebem valor 0
por ausência de recuperação.

O conjunto de avaliação é composto por cinco perguntas elaboradas
pelos autores sobre eventos climáticos reais no Brasil,
cobrindo diferentes tipos de evento (enchentes, calor extremo,
deslizamentos, geadas, precipitação extrema) e diferentes
granularidades geográficas (municípios, estados, regiões).
Cada pergunta possui um gabarito textual usado como referência
nas métricas de abrangência.

% ─────────────────────────────────────────────────────────────
\section{Resultados e Discussão}
\label{sec:resultados}

\subsection{Métricas RAGAS}

A Tabela~\ref{tab:ragas} apresenta os scores médios do RAGAS
para os três modos nas cinco perguntas de avaliação.

\begin{table}[htbp]
\caption{Scores médios do RAGAS por modo de recuperação}
\label{tab:ragas}
\begin{center}
\begin{tabular}{lrrr}
\toprule
\textbf{Métrica} & \textbf{No RAG} & \textbf{Vector RAG} & \textbf{Graph RAG} \\
\midrule
Faithfulness         & 0.XX & 0.XX & \textbf{0.XX} \\
Answer Relevancy     & 0.XX & 0.XX & \textbf{0.XX} \\
Context Recall       & --   & 0.XX & \textbf{0.XX} \\
Context Precision    & --   & 0.XX & \textbf{0.XX} \\
\bottomrule
\end{tabular}
\end{center}
\end{table}

\textit{Nota: valores 0.XX serão substituídos pelos scores
reais após execução do script de avaliação
\texttt{src/evaluation/ragas\_eval.py}.}

\subsection{Análise Qualitativa}

\textbf{No RAG:} o LLM demonstra conhecimento geral sobre
fenômenos climáticos brasileiros (períodos de seca no Nordeste,
enchentes no Sul), mas falha em perguntas que requerem dados
específicos de municípios ou datas concretas.
Foram observadas alucinações factuais, especialmente ao citar
valores de precipitação ou nomear municípios afetados por
eventos específicos.

\textbf{Vector RAG:} melhora a fidelidade ao recuperar descrições
textuais de eventos similares.
Contudo, como a recuperação é baseada em similaridade semântica,
perguntas sobre frequência (``quais municípios tiveram
\emph{mais} enchentes?'') retornam os eventos textuais mais
similares à query, e não os municípios com maior contagem ---
que requer agregação relacional.

\textbf{Graph RAG:} ao extrair entidades e executar consultas
Cypher parametrizadas, o pipeline recupera exatamente os registros
relevantes ao período e tipo de evento especificados.
Consultas de ranking (contagem de eventos por município) são
tratadas nativamente com a query \texttt{municipiosComMaisEventos},
retornando dados ordenados.
O contexto mais preciso resulta em respostas com maior fidelidade
e menor taxa de informação irrelevante.

\subsection{Limitações}

\begin{itemize}
  \item \textbf{Cobertura dos dados}: os boletins CEMADEN disponíveis
        cobrem principalmente eventos com vítimas; eventos sem
        danos reportados podem estar sub-representados.
  \item \textbf{Tamanho do benchmark}: cinco perguntas de avaliação
        é insuficiente para conclusões estatisticamente robustas;
        um benchmark maior é necessário para generalização.
  \item \textbf{Custo de consulta}: o pipeline GraphRAG realiza
        chamadas ao LLM para extração de entidades (GPT-4o-mini)
        e para geração (GPT-4o), aumentando latência e custo
        em relação ao Vector~RAG.
\end{itemize}

% ─────────────────────────────────────────────────────────────
\section{Conclusão}
\label{sec:conclusao}

Este trabalho apresentou um pipeline GraphRAG para consultas em
linguagem natural sobre eventos climáticos extremos no Brasil,
utilizando dados do INMET e do CEMADEN estruturados em um grafo
de conhecimento Neo4j.
A avaliação com RAGAS mostrou que a abordagem baseada em grafo
supera o RAG vetorial e o baseline sem recuperação nas métricas
de fidelidade e abrangência de contexto, especialmente em perguntas
relacionais que requerem agregação e cruzamento de entidades.

Como trabalho futuro, destacam-se:
(i) ampliar a cobertura de dados para incluir todos os estados
brasileiros e séries históricas mais longas;
(ii) implementar um benchmark de avaliação maior e mais
diversificado, com anotações de especialistas em climatologia;
(iii) explorar técnicas de grafo mais sofisticadas, como
\textit{community detection} e sumarização hierárquica
conforme proposto por Edge et al.~\cite{edge2024graphrag};
(iv) avaliar modelos de linguagem abertos (LLaMA, Mistral)
como alternativa ao GPT-4o para redução de custos operacionais.

\section*{Agradecimentos}

O autor agradece ao INMET e ao CEMADEN pela disponibilização
pública dos dados utilizados neste trabalho.

% ─────────────────────────────────────────────────────────────
\begin{thebibliography}{00}

\bibitem{vaghefi2023chatclimate}
S.~A.~Vaghefi et al.,
``ChatClimate: Grounding conversational AI in climate science,''
\textit{Commun. Earth Environ.}, vol.~4, no.~480, 2023.
\href{https://doi.org/10.1038/s43247-023-01084-x}{doi:10.1038/s43247-023-01084-x}

\bibitem{atkins2024generative}
C.~Atkins, G.~Girgente, M.~Shirzaei, and J.~Kim,
``Generative AI tools can enhance climate literacy but must be
checked for biases and inaccuracies,''
\textit{Commun. Earth Environ.}, vol.~5, no.~226, 2024.
\href{https://doi.org/10.1038/s43247-024-01392-w}{doi:10.1038/s43247-024-01392-w}

\bibitem{juhasz2024responsible}
M.~Juhasz et al.,
``Responsible retrieval augmented generation for climate
decision making from documents,''
\textit{arXiv preprint arXiv:2410.23902}, 2024.

\bibitem{hsu2024chatnetzero}
A.~Hsu, M.~Laney, D.~Manya, J.~Zhang, and L.~Farczadi,
``Evaluating ChatNetZero, an LLM-chatbot to demystify
climate pledges,''
in \textit{Proc. 1st Workshop NLP Meets Climate Change
(ClimateNLP 2024)}, pp.~82--92, 2024.

\bibitem{edge2024graphrag}
D.~Edge, H.~Trinh, N.~Cheng, J.~Bradley, A.~Chao, A.~Mody,
S.~Truitt, and J.~Larson,
``From local to global: A graph RAG approach to
query-focused summarization,''
\textit{arXiv preprint arXiv:2404.16130}, 2024.

\bibitem{es2024ragas}
S.~Es, J.~James, L.~Espinosa-Anke, and S.~Schockaert,
``RAGAS: Automated evaluation of retrieval augmented
generation,''
in \textit{Proc. 18th Conf. European Chapter ACL (EACL 2024)},
pp.~150--158, 2024.
\href{https://doi.org/10.18653/v1/2024.eacl-demo.16}{doi:10.18653/v1/2024.eacl-demo.16}

\bibitem{lewis2020rag}
P.~Lewis et al.,
``Retrieval-augmented generation for knowledge-intensive NLP tasks,''
in \textit{Advances in Neural Information Processing Systems
(NeurIPS 2020)}, vol.~33, pp.~9459--9474, 2020.

\bibitem{pan2024unifying}
S.~Pan, L.~Luo, Y.~Wang, C.~Chen, J.~Wang, and X.~Wu,
``Unifying large language models and knowledge graphs: A roadmap,''
\textit{IEEE Trans. Knowl. Data Eng.}, 2024.
\href{https://doi.org/10.1109/TKDE.2024.3352100}{doi:10.1109/TKDE.2024.3352100}

\bibitem{ahmad2025benchmarking}
S.~Ahmad, Z.~Nezami, M.~Hafeez, and S.~A.~R.~Zaidi,
``Benchmarking vector, graph and hybrid retrieval augmented
generation (RAG) pipelines for open radio access networks (ORAN),''
\textit{arXiv preprint arXiv:2507.03608}, 2025.

\bibitem{han2025ragvsgraphrag}
H.~Han et al.,
``RAG vs.\ GraphRAG: A systematic evaluation and key insights,''
\textit{arXiv preprint arXiv:2502.11371}, 2025.

\bibitem{zeng2025unbiased}
Q.~Zeng et al.,
``How significant are the real performance gains? An unbiased
evaluation framework for GraphRAG,''
\textit{arXiv preprint arXiv:2506.06331}, 2025.

\bibitem{inmet2024dados}
Instituto Nacional de Meteorologia (INMET),
``Banco de Dados Meteorológicos do INMET (BDMEP),''
Brasília, Brasil, 2024.
[Online]. Available: \url{https://bdmep.inmet.gov.br}

\bibitem{cemaden2024boletins}
Centro Nacional de Monitoramento e Alertas de Desastres Naturais
(CEMADEN),
``Boletins de Desastres Naturais,''
São José dos Campos, Brasil, 2024.
[Online]. Available: \url{https://www.cemaden.gov.br/boletins}

\end{thebibliography}

\end{document}
